% =============================================================================
\section{Kernel Modules}
\label{sec:kernelModules}
% =============================================================================
% -----------------------------------------------------------------------------
\subsection{Two- and Three-Dimensional Kernel Bindings}
\label{ssec:modules:kernel}
% -----------------------------------------------------------------------------
The \iiDandiiiDGeometryKernelPackage{}
package~\cite{cgal:bfghhkps-lgk23-21a} of \cgal{} consists of constant-size non-modifiable geometric primitive objects and operations on these objects. The objects are sets of points in
$d$-dimensional affine Euclidean space, where $d = 2,3$.  Each point is uniquely represented either by Cartesian coordinates or by homogeneous coordinates. An object type can be defined either precatively as a member of a kernel type or imperatively as a global class-template parameterized by a kernel type, defined in the C++ \cppLstinline{CGAL} namespace. For example, assume that \cppLstinline{Kernel} is a kernel type; the type that represents a two-dimensional point of this kernel is either \cppLstinline{Kernel::Point_2} or \cppLstinline{CGAL::Point_2<Kernel>}.  We generate bindings only for the latter types; that is, defining a two-dimensional point in Python amounts to the code below.\footnote{When writing code in C++ the precative style is advantageous, as it enables the extension of kernel object types, which is irrelevant when generating bindings.}
\begin{lstlisting}[style=pythonStyle]
>>> Ker = CGALPY.Ker      # define the module namespace
>>> Point_2 = Ker.Point_2 # define the point type in the module namespace
>>> p = Point_2(0, 0);
\end{lstlisting}
Table~\ref{tab:ker} lists the \cmake{} flags associated with the
\iiDandiiiDGeometryKernelPackage{} module.  The kernel type is an
instance of a chain of C++ class templates. For convenience, \cgal{}
provides the following predefned types of generally useful kernels:
\begin{compactenum}
\item \cppLstinline{Exact\_predicates\_inexact\_constructions\_kernel}---provides exact geometric predicates, but geometric constructions may be inexact due to round-off errors.
  %
\item \cppLstinline{Exact\_predicates\_exact\_constructions\_kernel}---provides exact geometric constructions, in addition to exact geometric predicates.
  %
\item \cppLstinline{Exact\_predicates\_exact\_constructions\_kernel\_with\_sqrt}---same as
  \cppLstinline{Exact\_predicates\_exact\_constructions\_kernel}, but the
  number type is a model of the concept that requires operations that perform
  square roots, namely \cppLstinline{FieldWithSqrt}.\footnote{See \url{https://doc.cgal.org/latest/Algebraic_foundations/classFieldWithSqrt.html}.}
  %
\item \cppLstinline{Exact\_predicates\_exact\_constructions\_kernel\_with\_kth\_root}---same as \cppLstinline{Exact\_predicates\_exact\_constructions\_kernel}, but the number type is a model of the concept that requires operations that perform $k$-th roots, namely   \cppLstinline{FieldWithKthRoot}.\footnote{See \url{https://doc.cgal.org/latest/Algebraic_foundations/classFieldWithKthRoot.html}}
  %
\item \cppLstinline{Exact\_predicates\_exact\_constructions\_kernel\_with\_root\_of}---same as \cppLstinline{Exact\_predicates\_exact\_constructions\_kernel}, but the number type is a model of the concept that requires operations that computes the root of univariate polynomial, namely \cppLstinline{FieldWithRootOf}.\footnote{See \url{https://doc.cgal.org/latest/Algebraic_foundations/classFieldWithRootOf.html}.}
\end{compactenum}

%%%%%%%% Kernel Module Flags
\begin{table}[!htp]
  \caption{\iiDandiiiDGeometryKernelPackage{} module flags}
  \centering\begin{tikzpicture}
  \matrix (first) [tableModuleFlags,
    row 3/.style={nodes={minimum height=2.8em}}
  ] {
    Name & Type & Default & Description\\
    \cmakeLstinline{KERNEL\_NAME} & String & \myLstinline{epic} &
      The kernel type used\\
    \cmakeLstinline{KERNEL\_INTERSECTION\_BINDINGS} & Boolean &
      \myLstinline{true} &
      Determines whether to generate bindings for intersections\\
  };
 \end{tikzpicture}
 \label{tab:ker}
\end{table}

All the predefined types are of Cartesian kernels. The \cmake{} flag \myLstinline{KERNEL\_NAME} specifies which kernel type should be used for the generated bindings.  Currently, only a limited predefined types and two specific types are supported; see Table~\ref{tab:ker:name}.

%%%%%%%% Kernel Name Options
\begin{table}[!htp]
  \caption{Kernel name options. \cppLstinline{pre..} stands for
    \cppLstinline{predicates}.}
 \centering\begin{tikzpicture}
   \matrix (first) [tableName,
     row 6/.append style={nodes={minimum height=2.8em}},
     column 1/.style={nodes={text width=19em}},
     column 2/.style={nodes={text width=30em}}
   ] {
   \cmakeLstinline{KERNEL\_NAME} & Predefined Type\\
   \myLstinline{epic} &
     \cppLstinline{Exact\_pre...\_inexact\_constructions\_kernel}\\
   \myLstinline{epec} &
     \cppLstinline{Exact\_pre...\_exact\_constructions\_kernel}\\
   \myLstinline{epecws} &
     \cppLstinline{Exact\_pre...\_exact\_constructions\_kernel\_with\_sqrt}\\
   \myLstinline{filteredSimpleCartesianDouble} &
     \cppLstinline{Filtered\_kernel<Simple\_cartesian<double>>}\\
   \myLstinline{filteredSimpleCartesianLazyGmpq} &
     \setlength{\tabcolsep}{0pt}\renewcommand{\arraystretch}{0}
     \begin{tabular}{l}
       \cppLstinline{NT = Lazy\_exact\_nt<Gmpq>}\\[5pt]
       \cppLstinline{Filtered\_kernel<Simple\_cartesian<NT>>}
     \end{tabular}\\
  };
 \end{tikzpicture}
 \label{tab:ker:name}
\end{table}

The kernel type determines the underlying number type used to represent coefficients and coordinates of kernel objects and for evaluating mathematical expressions that involve these coefficients and coordinates. The Python attribute \pyLstinline{FT}, nested in the Python namespace \pyLstinline{Ker}, exposes the C++ underlying number type when it is not a primitive data type; that is, when the selected kernel is nor \cppLstinline{Exact\_predicates\_inexact\_constructions\_kernel} neither \cppLstinline{Filtered\_kernel<Simple\_cartesian<double>>}. (In both cases the underlying number type is \cppLstinline{double}.) Similar to The Python code above, the code below defines a
two-dimensional point; here, the coordinates are explicitly converted to the underlying number type.
\begin{lstlisting}[style=pythonStyle]
>>> p = Point_2(Ker.FT(0), Ker.FT(0))
\end{lstlisting}

The code excerpt below, in the \cmake{} language, sets the \cmake{} flags for our first example. When \myLstinline{cmake} is applied with these settings followed by the native build commands (e.g., \myLstinline{make} on Linux platforms) a library called \myLstinline{CGALPY} is generated. This library consists of the bindings necessary to run the Python example that follows. Observe
that bindings of the \iiDandiiiDGeometryKernelPackage{} module are generated by default, whereas bindings for all other modules are not generated by default.
\framedInputCmake[\normalsize]{../../../cmake/tests/release/epec_fixed_release.cmake}

The following code example in Python computes the intersection point of two segments using the generated binding.
\framedInputPython[\normalsize][14]{../../../src/python_scripts/cgalpy_examples/ker_intersection.py}

% -----------------------------------------------------------------------------
\subsection{$d$-Dimensional Kernel Bindings}
\label{ssec:modules:kerneld}
% -----------------------------------------------------------------------------
\cgal{} includes a separate package that consists of constant-size non-modifiable geometric primitive objects in arbitrary dimensions, and operations on these objects, called
\dDGeometryKernelPackage{}~\cite{cgal:s-gkd-21a}. Similar to the two- and three-dimensional kernels, the objects of the $d$-dimensional kernels are sets of points in some $d$-dimensional affine Euclidean space, where the dimension $d$ is either static across a kernel type
or dynamic; see below for more details. Each point is uniquely represented either by Cartesian coordinates or by homogeneous coordinates. For convenience, \cgal{} provides the following predefined types of generally useful kernels:
\begin{compactenum}
\item \cppLstinline{Epick\_d<DimensionTag>}---provides exact geometric
  predicates, but geometric constructions may be inexact due to
  round-off errors.
  %
\item \cppLstinline{Epeck\_d<DimensionTag>}---provides exact geometric
  constructions, in addition to exact geometric predicates.
  %
\end{compactenum}
Table~\ref{tab:kerd} lists the \cmake{} flags associated with the
\dDGeometryKernelPackage{} module.

%%%%%%%% D Kernel Module Flags
\begin{table}[!htp]
  \caption{\dDGeometryKernelPackage{} module flags}
  \centering\begin{tikzpicture}
  \matrix (first) [tableModuleFlags,
    column 3/.style={nodes={text width=4.5em}},
    column 4/.style={nodes={text width=20.8em}}
  ] {
    Name & Type & Default & Description\\
    \cmakeLstinline{KERNEL\_D\_NAME} & String & \myLstinline{epicd} &
      The kernel type used\\
    \cmakeLstinline{KERNEL\_D\_DIMENSION\_TAG} & String & dynamic &
      Determines whether the dimension is dynamic\\
    \cmakeLstinline{KERNEL\_D\_DIMENSION} & Integer & \myLstinline{2} &
      The dimension of the ambient Euclidean space\\
  };
 \end{tikzpicture}
 \label{tab:kerd}
\end{table}

The \cmake{} flag \cmakeLstinline{KERNEL\_D\_NAME} specifies which
kernel type should be used for the generated bindings; see
Table~\ref{tab:kerd:name}.

%%%%%%%% D Kernel Name Options
\begin{table}[!htp]
  \caption{$d$-dimensional Kernel name options.}
  \centering\begin{tikzpicture}
  \matrix (first) [tableName, column
    1/.style={nodes={text width=19em}},
    column 2/.style={nodes={text width=30em}}
  ] { \cmakeLstinline{KERNEL\_D\_NAME} & Predefined Type\\
    \cmakeLstinline{epicd} & \cppLstinline{Epick\_d<DimensionTag>}\\
    \cmakeLstinline{epecd} & \cppLstinline{Epeck\_d<DimensionTag>}\\
    \cmakeLstinline{cartesiandDouble} & \cppLstinline{Cartesian_d<double>}\\
    \cmakeLstinline{cartesiandLazyGmpq} &
      \cppLstinline{Cartesian\_d<Lazy\_exact\_nt<Gmpq>>}\\
  };
 \end{tikzpicture}
 \label{tab:kerd:name}
\end{table}

When either \cppLstinline{Epick\_d<DimensionTag>} or \cppLstinline{Epeck\_d<DimensionTag>} are instantiated the template parameter must be substituted with a type that represents the dimension of the ambient Euclidean space. It may be either \cppLstinline{Dimension\_tag<d>} where $d$ is an integer or \cppLstinline{Dynamic\_dimension\_tag}. In the latter case, the dimension of the space is specified for each point when it is constructed, so it does not need to be known at compile-time of the bindings. The \cmake{KERNEL\_D\_DIMENSION\_TAG} flag specifies whether the dimension is static or dynamic. If it is static, the dimension is extracted from the \cmake{KERNEL\_D\_DIMENSION} \cmake{} flag; see Table~\ref{tab:kerd:dim}.

%%%%%%%% D Kernel Dimension Tag Options
\begin{table}[!htp]
  \caption{$d$-dimensional Kernel dimension tag options.}
  \centering\begin{tikzpicture}
  \matrix (first) [tableName, column
    1/.style={nodes={text width=19em}},
    column 2/.style={nodes={text width=30em}}
  ] { \cmakeLstinline{KERNEL\_D\_DIMENSION\_TAG} & Predefined Type\\
    \cmakeLstinline{static} & \cppLstinline{Dimension_tag<d>}\\
    \cmakeLstinline{dynamic} & \cppLstinline{Dynamic\_dimension\_tag}\\
  };
 \end{tikzpicture}
 \label{tab:kerd:dim}
\end{table}

The kernel type determines the underlying number type. It is possible to have different underlying number types for the \iiDandiiiDGeometryKernelPackage{} and the \dDGeometryKernelPackage{} models. However, when the number types differ, expensive conversions
might be necessary to combine operations from both kernels, or it may not be possible at all using the current bindings.

The code excerpt below, in the \cmake{} language, sets the \cmake{} flags for our second example. When \myLstinline{cmake} is applied with these settings followed by the native build commands a library called \myLstinline{CGALPY_kerdCdlgDynamic} is generated. This library consists of the bindings necessary to run the Python example that follows. \framedInputCmake[\normalsize]{../../../cmake/tests/release/cdlg_release.cmake}

The following code example in Python determines whether two segments in four dimensions intersect using the generated binding.
\framedInputPython[\normalsize][14]{../../../src/python_scripts/cgalpy_examples/kerd_do_intersect.py}
